\section{Evaluation}
\label{sec:evaluation}

\subsection{Research Questions}

We organize the evaluation around three questions:

\begin{description}
  \item[RQ1:] How effectively does \tool find previously unknown bugs in real
  Python native extensions?
  \item[RQ2:] How does \tool compare with existing Python native-extension bug
  detectors, and why does it miss bugs found by another detector?
  \item[RQ3:] How do static pattern-based pruning and structured prompt
  design affect the cost and bug recall of \tool?
\end{description}

RQ1 evaluates the end-to-end result after the high-confidence filter. RQ2
compares with the closest reproducible analyzer and studies non-overlapping
defects. RQ3 isolates the two design decisions that control analysis cost and
model behavior.

\subsection{Subjects and Protocol}

We selected nine open-source projects according to three criteria:
(1) the project implements substantial performance-critical or system-facing
functionality in C or C++; (2) the project is actively maintained with a
public issue tracker; and (3) the project is widely adopted, as reflected
by GitHub star count and PyPI download volume. Table~\ref{tab:subjects} reports the quantitative
corpus. The scanner manifest contains 489 C/C++ source files from these
projects.

\begin{table}[t]
\caption{Subjects and manually audited high-confidence PyCGuard reports.}
\label{tab:subjects}
\centering
\small
\setlength{\tabcolsep}{3.5pt}
\begin{tabular}{l|r|r|r|r|r}
\toprule
Project & Stars & Bug Reports & TP & FP & Precision \\
\midrule
SciPy      & 14.8k & 24 & 20 & 4 & 83.3\% \\
pandas     & 49.1k & 18 & 15 & 3 & 83.3\% \\
NumPy      & 32.3k & 15 & 13 & 2 & 86.7\% \\
Pillow     & 13.6k & 15 & 13 & 2 & 86.7\% \\
PyTorch    & 101k  & 12 & 10 & 2 & 83.3\% \\
PycURL     &  1.2k &  7 &  5 & 2 & 71.4\% \\
ujson      &  4.5k &  6 &  5 & 1 & 83.3\% \\
TensorFlow & 196k  &  4 &  3 & 1 & 75.0\% \\
psutil     & 11.2k &  3 &  2 & 1 & 66.7\% \\
\midrule
Total      &       & 104 & 86 & 18 & 82.7\% \\
\bottomrule
\end{tabular}
\end{table}

\paragraph{PyCGuard configuration.}
RQ1 uses the configuration in Section~\ref{sec:pipeline}, including
\code{claude-sonnet-4}, temperature 0, early stop, and the
high-confidence report filter. The analysis unit is one project, file,
function, and SO tuple. Several unsafe effects in one tuple are counted as one
report. We do not use lower-confidence violations to calculate precision.
Medium- and low-confidence outputs are suppressed and not audited;
all precision claims in this paper are restricted to the 104
high-confidence reports.

\paragraph{Manual audit.}
For each report, we inspect the complete function, the reported path, the
Python/C API contract, and relevant callees or macros when needed. A report is
a true positive when the reported behavior violates the claimed SO on a
feasible path. A report is a false positive when the path is infeasible, a
hidden ownership or cleanup action discharges the SO, or the model
misinterprets the API contract. Precision is $TP/(TP+FP)$. Because the study
starts from tool reports rather than a complete defect oracle, we do not claim
recall for the in-the-wild corpus.

\paragraph{Baseline protocol.}
We run PyRefcon over all C/C++ source files accepted by its compilation
pipeline. We parse 224 HTML warnings and group warnings with the same project,
file, function, and bug type into 80 audit units. Each unit is inspected using
the reported trace and source context. This deduplication prevents repeated
paths for one defect from inflating the result.

\subsection{RQ1: In-the-Wild Bug Finding}

\tool reports 104 high-confidence findings. Manual audit identifies 86 true
positives and 18 false positives, for a precision of
$86/104=82.7\%$. The true positives cover a diverse range of defect
classes, including reference ownership errors, null dereferences,
exception-state violations, resource leaks, post-release safety errors,
type mismatches, and numeric bounds violations. This distribution
supports the central scope claim: Python native-extension defects extend
beyond reference-count imbalance.

The project-level result is not dominated by one code base. Seven projects
contribute at least five true positives, and all nine projects in the audited
corpus contain at least two. Of these findings, 45 were accepted by project
maintainers, of which 13 have been fixed and 32 are confirmed but awaiting
a patch.

The motivating defect in Figure~\ref{fig:motivation} illustrates a common
benefit of obligation checking. A shared cleanup block appears locally
responsible because it decrements every array variable. The ownership
obligation forces the checker to distinguish the borrowed and new-reference
branches. Similar true positives require relating a fallible allocation to a
later unsafe macro, preserving exception state across callbacks, or pairing a
buffer acquisition with all failure exits.

\paragraph{False positives.}
Table~\ref{tab:fp} categorizes the 18 false positives by their immediate root
cause. Six require precise knowledge of whether an API borrows, creates, or
steals a reference. The next largest group requires reasoning about complex
control flow. Macros, C++ destructors, and paired teardown functions hide
discharge evidence outside the source form given to the model. These failures
suggest that the dominant limitation is missing semantic context rather than a
lack of obligation categories.

\begin{table}[t]
\caption{Root causes of the 18 false-positive PyCGuard reports.}
\label{tab:fp}
\centering
\small
\begin{tabular}{lr}
\toprule
Root cause & Count \\
\midrule
Borrowed, new, or stolen reference semantics & 6 \\
Complex control flow & 3 \\
Macro expansion & 2 \\
C++ RAII or smart-pointer cleanup & 2 \\
Interprocedural cleanup & 2 \\
State-machine behavior & 1 \\
Standard error propagation & 1 \\
Infeasible retry path & 1 \\
\bottomrule
\end{tabular}
\end{table}

\subsection{RQ2: Comparison with Existing Tools}

PyRefcon is the closest reproducible detector because it directly analyzes
Python native code and models Python object lifecycle and reference counts
\cite{pyrefcon2023}. We run PyRefcon over all 489 C/C++ source files in
the nine projects and manually audit each deduplicated warning unit.
Table~\ref{tab:baseline} reports the per-project result.

\begin{table}[t]
\caption{PyRefcon per-project audit result on the nine subject projects.}
\label{tab:baseline}
\centering
\small
\begin{tabularx}{\columnwidth}{X|r|r|r|r|r}
\toprule
Project & Warnings & Units & TP & FP & Precision \\
\midrule
SciPy      & 76 & 11 &  1 & 10 & 9.1\%  \\
pandas     & 32 & 16 &  4 & 12 & 25.0\% \\
NumPy      & 34 & 14 &  1 & 13 & 7.1\%  \\
Pillow     & 34 & 17 &  0 & 17 & 0.0\%  \\
PyTorch    & \texttimes & \texttimes & \texttimes & \texttimes & \texttimes \\
PycURL     & 26 & 12 &  0 & 12 & 0.0\%  \\
ujson      &  6 &  3 &  0 &  3 & 0.0\%  \\
TensorFlow &  6 &  2 &  0 &  2 & 0.0\%  \\
psutil     & 10 &  5 &  0 &  5 & 0.0\%  \\
\midrule
Total      & 224 & 80 &  6 & 74 & 7.5\%  \\
\bottomrule
\end{tabularx}
\end{table}

PyRefcon emits 224 raw warnings across 38 of 489 target files; 151 of
the remaining files fail to compile in its analysis environment.
PyTorch fails entirely because its required CMake-generated headers
are absent from the standalone analysis toolchain.
After deduplication, 80 audit units remain, of which only 6 are true
positives (7.5\% precision).

The low true-positive count reflects a fundamental scope mismatch.
PyRefcon models only reference-count ownership, so it cannot represent
exception-state rules, resource-pairing obligations, type-agreement
requirements, numeric bounds, or post-release safety. These five categories
account for 35 of PyCGuard's 86 true positives and are entirely outside
PyRefcon's scope. Even within
reference counting, one of the six PyRefcon true positives overlaps a
PyCGuard finding; the remaining five are raw-pointer null dereferences
or generic \code{malloc}/\code{free} flows whose correctness depends on
project-local C invariants rather than a Python/C API contract, so
PyCGuard's static binder never creates an obligation for them.

The 74 false positives reveal the cost of modeling a single lifecycle
pattern without explicit discharge alternatives. Fifty-one reports trace
a standard \code{Py\_DECREF} call into the CPython deallocation body
and flag it as an ownership release error, treating any decrement as
suspect regardless of whether the preceding acquire-release pair is
balanced. Eight misread cleanup paths that explicitly destroy an
already-allocated object. Nine miss cross-function or cross-iteration
invariants, and six arise from macro expansion that hides the actual
reference semantics. These patterns are precisely the failure modes
that PyCGuard's discharge-condition list is designed to rule out: each
discharge condition represents a syntactic form that, when present on
every exit path, provides evidence that the obligation is satisfied without requiring
interprocedural or macro-expansion reasoning.

We attempted to run CpyChecker using the configuration reported by
PyRefcon. The plugin either disabled reference-count verification on
GCC 7 and later, failed to initialize with the matching Python headers,
or produced no plugin warnings in the usable configuration. We
therefore report it as unavailable. Pungi has no public implementation.
RID targets general reference-count inconsistencies on the Linux kernel
\cite{rid2016} and is not a direct Python/C baseline. CHARON analyzes
cross-language vulnerability flows \cite{charon2025} with no public
artifact for a same-project run. These tools are discussed qualitatively
and do not receive quantitative rows.

\subsection{RQ3: Design Analysis}

RQ3 quantifies the value of two design decisions that govern the
end-to-end behavior of \tool. The first is static pattern-based
pruning, which decides how many obligations reach the LLM. The second
is the obligation prompt, which decides how the LLM reasons about each
obligation that it does see. We measure pruning by its marginal effect
on checking cost over the 86 manually validated true-positive cases;
this controlled replay holds the analyzed function/SO set fixed and
therefore isolates the cost contribution of static pruning. We measure
the prompt by its effect on bug recall across the same 86 cases.

\paragraph{Cost effect of pruning.}
Table~\ref{tab:pruning-cost} reports per-project token usage, dollar
cost, and wall-clock time for a controlled replay of the 86 RQ1
true-positive cases with pruning disabled and enabled. The same
project/file/function/SO targets are used in both configurations;
differences in token usage and time are therefore attributable to
pruning rather than to different report populations. Cost uses
Anthropic's published Claude Sonnet 4 pricing
(\$3 / MTok input, \$15 / MTok output). The high-confidence filter and
the early-stop rule from Section~\ref{sec:pipeline} are kept enabled
in both configurations so that the comparison reflects only the
pruning step.

\begin{table*}[t]
\caption{Per-project LLM cost on the 86 true positives with and without static pruning.}
\label{tab:pruning-cost}
\centering
\small
\setlength{\tabcolsep}{4.5pt}
\begin{tabular}{l rrrr c rrrr c r}
\toprule
& \multicolumn{4}{c}{Without pruning} & & \multicolumn{4}{c}{With pruning} & & \\
\cmidrule{2-5} \cmidrule{7-10}
Project & Input (K) & Output (K) & Cost (\$) & Time (min) & &
          Input (K) & Output (K) & Cost (\$) & Time (min) & & Cost Reduction (\%) \\
\midrule
SciPy      & 103.0 & 24.7 & 0.68 & 8.0 & &  62.3 & 17.1 & 0.44 & 5.5 & & 35.3\% \\
pandas     &  28.4 & 13.8 & 0.29 & 4.4 & &  17.6 &  9.8 & 0.20 & 3.1 & & 31.0\% \\
NumPy      &  81.6 & 15.7 & 0.48 & 5.3 & &  47.5 & 10.4 & 0.30 & 3.5 & & 37.5\% \\
Pillow     &  37.9 & 15.0 & 0.34 & 4.8 & &  22.7 & 10.2 & 0.22 & 3.2 & & 35.3\% \\
PyTorch    &  15.5 &  8.4 & 0.17 & 2.7 & &  11.9 &  7.4 & 0.15 & 2.4 & & 11.8\% \\
PycURL     &  10.8 &  7.4 & 0.14 & 2.3 & &   6.0 &  4.6 & 0.09 & 1.4 & & 35.7\% \\
ujson      &   6.3 &  4.7 & 0.09 & 1.6 & &   4.1 &  3.5 & 0.06 & 1.2 & & 33.3\% \\
TensorFlow &   3.4 &  2.7 & 0.05 & 0.9 & &   2.4 &  2.2 & 0.04 & 0.7 & & 20.0\% \\
psutil     &   9.3 &  2.4 & 0.06 & 0.8 & &   4.8 &  1.4 & 0.04 & 0.4 & & 33.3\% \\
\midrule
Total      & 296.2 & 94.9 & 2.31 & 30.7 & & 179.2 & 66.7 & 1.54 & 21.5 & & 33.3\% \\
\bottomrule
\end{tabular}
\end{table*}

Processing the 86 true-positive cases with pruning disabled costs
\$2.31 and takes 30.7 minutes of wall-clock LLM time. Enabling
pruning cuts this to \$1.54 and 21.5 minutes, a reduction of 33.3\%
in cost and 30\% in time. The driver is a sharp drop at the
binding stage: pruning eliminates a substantial fraction of bound
obligations per project before any model call, so
the LLM never pays the per-call overhead for cases it would
ultimately answer SAFE. The end-to-end reduction is smaller than the
binding-stage reduction because the obligations that survive pruning
are harder on average; the model emits longer structured rationale on them,
which cost more output tokens per call. Pruning therefore
concentrates LLM effort on syntactically suspicious cases rather than
removing work uniformly.

Per-project savings track the fraction of callback-heavy and
allocator-only functions that the binders hit but pruning can
dismiss. SciPy, NumPy, PycURL, and psutil each achieve at least 33\%
cost reduction, the highest cuts occurring where short thunk or
conversion helpers account for most of the unpruned obligation count.
PyTorch shows the smallest reduction (11.8\%) because its
mixed-language RAII helpers survive pruning at a higher rate; the
surviving obligations are in larger, more complex functions, and the
structured rationale on them is longer. Even there, every audited true
positive is retained, confirming that the risk indicator set
does not discard evidence for any of the audited true-positive cases.

\paragraph{Bug recovery effect of CoT prompting and Discharge
Conditions.}
We next isolate the two reasoning-side design decisions: the
chain-of-thought (CoT) discharge procedure encoded in the system
prompt, and the per-obligation discharge-condition (DC) list provided
in the user prompt. We rerun the 86 RQ1 true positives under two
ablated configurations and otherwise hold the pipeline fixed.
\textit{w/o CoT} replaces the CoT system prompt with a direct-verdict
instruction. \textit{w/o DC} retains the CoT system prompt but removes
the DC list from the user prompt; the obligation text alone tells the
model what to discharge. The baseline for comparison is the RQ1
configuration (CoT prompt with DC list). The backend, temperature, max
tokens, high-confidence filter, and obligation binder are identical
across all configurations. We classify a case as \emph{recovered} when
the model returns parseable JSON with \code{verdict=VIOLATED} and
\code{confidence=high}. This ablation is not a precision experiment:
because it is run only on independently audited true positives, it
measures how often each prompt variant preserves known bug reports; it
does not characterize false-positive behavior or deployment precision.

\begin{table}[t]
\caption{Per-project bug recall on the 86 true positives under two ablations.}
\label{tab:llm-ablation}
\centering
\small
\setlength{\tabcolsep}{4pt}
\begin{tabular}{lr rr rr}
\toprule
& & \multicolumn{2}{c}{w/o CoT} & \multicolumn{2}{c}{w/o DC} \\
\cmidrule(lr){3-4} \cmidrule(lr){5-6}
Project    & Target & Found & Recall & Found & Recall \\
\midrule
SciPy      & 20 & 17~{\scriptsize\textcolor{red}{($\downarrow$3)}} & 85.0\% & 19~{\scriptsize\textcolor{red}{($\downarrow$1)}} & 95.0\% \\
pandas     & 15 & 15~{\scriptsize\textcolor{green!60!black}{($\downarrow$0)}} & 100.0\% & 15~{\scriptsize\textcolor{green!60!black}{($\downarrow$0)}} & 100.0\% \\
NumPy      & 13 &  5~{\scriptsize\textcolor{red}{($\downarrow$8)}} & 38.5\% & 10~{\scriptsize\textcolor{red}{($\downarrow$3)}} & 76.9\% \\
Pillow     & 13 & 10~{\scriptsize\textcolor{red}{($\downarrow$3)}} & 76.9\% & 10~{\scriptsize\textcolor{red}{($\downarrow$3)}} & 76.9\% \\
PyTorch    & 10 &  9~{\scriptsize\textcolor{red}{($\downarrow$1)}} & 90.0\% & 10~{\scriptsize\textcolor{green!60!black}{($\downarrow$0)}} & 100.0\% \\
PycURL     &  5 &  2~{\scriptsize\textcolor{red}{($\downarrow$3)}} & 40.0\% &  1~{\scriptsize\textcolor{red}{($\downarrow$4)}} & 20.0\% \\
ujson      &  5 &  3~{\scriptsize\textcolor{red}{($\downarrow$2)}} & 60.0\% &  3~{\scriptsize\textcolor{red}{($\downarrow$2)}} & 60.0\% \\
TensorFlow &  3 &  1~{\scriptsize\textcolor{red}{($\downarrow$2)}} & 33.3\% &  3~{\scriptsize\textcolor{green!60!black}{($\downarrow$0)}} & 100.0\% \\
psutil     &  2 &  2~{\scriptsize\textcolor{green!60!black}{($\downarrow$0)}} & 100.0\% &  2~{\scriptsize\textcolor{green!60!black}{($\downarrow$0)}} & 100.0\% \\
\midrule
Total      & 86 & 64~{\scriptsize\textcolor{red}{($\downarrow$22)}} & 74.4\% & 73~{\scriptsize\textcolor{red}{($\downarrow$13)}} & 84.9\% \\
\bottomrule
\end{tabular}
\end{table}

Table~\ref{tab:llm-ablation} reports the per-project recovery counts.
The ablation evaluates each of the 86 true-positive cases as an
independent unit under two ablated configurations and otherwise holds
the pipeline fixed.

The CoT step is the primary driver of recall. Without it, the tool
finds only 64 of the 86 known bugs, missing 22. These losses
concentrate on defect patterns whose discharge requires enumerating
multiple exit paths: 12 are reference-count leaks, borrowed-reference
errors, and resource leaks in NumPy, SciPy, PycURL, and Pillow where
the defect lies on an error-goto or non-normal-return branch that the
direct prompt does not enumerate; 2 are exception-state violations in
TensorFlow where an unchecked call on an alternate branch silently
clears the error indicator. Without the per-path enumeration that CoT
enforces, the model inspects only the dominant path and returns SAFE.

The DC list sharpens the model's judgment on individual obligations.
Without it, the tool finds 73 of the 86 known bugs, missing 13. These
losses concentrate on cases where the obligation description alone
leaves the violation criterion underspecified: resource leaks and
unchecked-return patterns where the model, lacking explicit discharge
alternatives to rule out, defaults to a permissive SAFE judgment. The
DC list narrows the reasoning to the specific failure pattern and
prevents this under-specification.

The two components address complementary failure modes: CoT prevents
path-coverage failures and the DC list prevents
obligation-underspecification failures. Removing either degrades
true-positive recovery on the 86-case set, with CoT contributing the
larger share.


\paragraph{Answers to the research questions.}
For RQ1, \tool finds 86 previously unknown bugs among 104 audited
high-confidence reports, with 82.7\% precision. For RQ2, PyRefcon
finds six true positives among 80 deduplicated units; five
non-overlapping findings are generic C memory errors outside the
Python/C API obligation scope. For RQ3, static pattern-based pruning
cuts the LLM cost by 33.3\% and wall-clock time by 30\% on the 86
true-positive cases while preserving every audited finding, and the
prompt ablation shows that removing the CoT step drops recall from
84.9\% to 74.4\% while removing the discharge-condition list shifts
the recovered set without changing the total recall.
